\documentclass[10pt]{article}
\usepackage[margin=20mm]{geometry}
\usepackage{amsmath,amssymb,fontspec,unicode-math,xurl,hyperref}
\IfFontExistsTF{DejaVu Serif}{\setmainfont{DejaVu Serif}}{\setmainfont{Latin Modern Roman}}
\setmathfont{latinmodern-math.otf}
\setlength{\parindent}{0pt}\setlength{\parskip}{6pt}\sloppy
\begin{document}
ELVIA Scientific Verification Report
0096 / I25-43465\par
From Mechanistic Interpretability to Mechanistic Biology: Training, Evaluating, and Interpreting Sparse Autoencoders on Protein Language Models\par
Reviewed version: PMLR 267 (ICML 2025), pp. 460-476, publisher-linked 17-page PDF. Review date: 24 October 2026. PDF physical pages 1-17 equal printed article pages 1-17; proceedings pagination is not used below.\par
Source SHA-256: cf3326a7c0dcfce98a6e66c5fd63891e2ebaaaf93e1f963dd8329f169cfedc37. Hash and title match the source-acquisition record; original bytes preserved.\par
Outcome: No Issue Found within the inspected scientific scope, with scope clarifications and unresolved evidence questions. Public implementation discrepancies are recorded separately; no demonstrated central-claim damage. Scientific decisions remain pending.\par
\section*{1. Scope and decisive assessment}
The inspected contribution is a method for extracting and characterizing ESM-2 SAE latents, comparing their interpretability with ESM coordinates, and using linear probes to nominate biologically meaningful features. The abstract and discussion explicitly frame unknown mechanisms as hypotheses for subsequent experimental testing. The ATP-binding/CHO-expression suggestion is not presented as an experimentally established causal mechanism.\par
Main text pp. 1-9 and technical appendix pp. 12-17 were read. Figures 4, 5 and 7 were visually inspected from the preserved PDF. Selected publisher-linked repository files were inspected statically at current commit fae2d3ae90220781289e27a26cb3b0826cf2b07e (23 April 2026) and publication-period commit 287c373817b06178603fc3a97bec95ebabb17ecc (19 May 2025). Availability of a commit does not establish that it generated the paper’s results.\par
No proof or empirical result was falsified by the checks below. This outcome is bounded by the coverage ledger; it is not certification of the whole paper.\par
\section*{2. Definitions, dependencies and checked reasoning}
Claim and location | Check and result\par
SAE and sparsity; Sec. 3.1, p. 3 | Checked encoder/decoder definitions against released model and training code. Top-k followed by ReLU yields at most k active latents. Release also normalizes inputs, uses an encoder bias and an auxiliary loss; simplified paper equations are insufficient to reconstruct every implementation detail, but this does not demonstrate failure of the learned features.\par
Family specificity; Secs. 4.2-4.4, p. 5; A.6, p. 13 | The operational label is F1 > 0.7 for a latent predicting a selected InterPro family. The inspected code tunes the threshold on training rows and evaluates F1 on test rows. Candidate families and normalization maxima are selected before this split; the stronger interpretation of fully independent target-family selection is not established.\par
Probe performance; Secs. 3.2-3.3 and 5.1; Fig. 5 | Read task metrics and stated homology-based splits; checked plotted direction of comparisons. Subcellular demo separates train/test and uses validation flags with PredefinedSplit and GridSearchCV. No test-label hyperparameter tuning was found in that path. Published scores and homology clustering were not reproduced.\par
Biological interpretation; Sec. 5.2, pp. 6-8; discussion | Specific motif associations are illustrative and the layer selection is explicitly qualitative. Unknown functional associations are hypotheses. Their causal validation would require fresh experiments; absence of those experiments does not refute the stated hypothesis-generation contribution.\par
\section*{3. Scoped unresolved questions and counterevidence}
Human ratings (Sec. 3.4, p. 4; Fig. 4, p. 5). Seven participants are reported; the figure explicitly uses binomial one-versus-all confidence intervals and shows about 80\% SAE ’yes’ ratings. Original rater/feature assignments and exports are needed to assess dependence and interval calibration. Dependence is not assumed to have a particular magnitude, and neither the descriptive fraction nor its direction was disproved. A publication-period UI advertises 100 features per model, whereas the text says 100 total, and uses randomized sorting rather than a uniform shuffle. These source discrepancies do not establish the actual deployed protocol or numerical effect.\par
Steering (Sec. 4.2, p. 5; A.5/Fig. 7, p. 13). The plot supports the reported descriptive ordering. With the other latents fixed, a single latent clamp changes the normalized decoded stream by the decoder direction times the latent increment. Equal clamp multipliers need not give equal increments or decoded perturbation norms. Thus the evolutionary-constraint explanation is suggestive; no measured ordering was overturned by this algebra. The article itself treats useful steering as an open direction.\par
Coefficient interpretation (A.10/Fig. 10, p. 17). The reported variance ratio is 0.00203 / 0.000324 = 6.2654. Coefficient dispersion alone does not identify why the features are more predictive or monosemantic. The secondary-structure training/preprocessing artifacts needed to adjudicate this explanatory statement were not located. A generic rescaling of coordinates is not asserted to be an admissible counterexample to the particular trained models.\par
Family evaluation (A.6). The released grid is 0.3-0.9 in seven steps, rather than the stated 0.1-0.9 in nine steps. This is a confirmed source-description discrepancy, including in the inspected publication-period file. Its effect on the reported family-specific counts is unknown. Selecting the candidate family from the highest-activating sequence is itself described in the paper; it must not be misreported as test-set threshold optimization.\par
\[\Delta x_{\mathrm{norm}}=d_j(a_j^{*}-z_j).\]\par
\section*{4. Existing public concern, versions and remaining work}
Repository issue \#44 (opened 4 December 2025) alleges that notebooks/sae\_test\_steering.ipynb skips later ESM layers. That target notebook was not acquired in this scoped selection, so the allegation was not independently reproduced. A different inspected path, interprot/steering.ipynb cell 8 and esm\_wrapper.py lines 88-99, explicitly resumes the remaining layers before the language-model head. This is counterevidence to treating every released steering path as faulty; it does not resolve issue \#44. No link from the alleged notebook behavior to Appendix A.5 experiment generation was established.\par
The bioRxiv record 10.1101/2025.02.06.636901v2 is dated 18 June 2025. Targeted publisher/preprint/repository searches identified no formal erratum relevant to the checked claims; this is a scoped search result. OpenReview returned a browser-verification wall, and direct bioRxiv opening returned 403. The publisher PDF remains the reviewed version.\par
Closing the open items needs original run commits/configurations, rating exports, dataset split identifiers, per-family outputs, steering experiment scripts and secondary-structure probe preprocessing. No author code, model inference, training, paid API request or new biological experiment was run. No scientifically verified repair is proposed.\par
\section*{5. Conclusion and decision separation}
Item | Issue judgment | Repair / adoption / remaining work\par
Core SAE/probing contribution | No Issue Found within checked scope | No repair required by this audit; scores and causal biology not reproduced.\par
Ratings, steering interpretation, coefficient explanation | Needs Evidence / scope clarification | Raw assignments, controlled perturbations and preprocessing needed. Pending questions do not raise supported severity.\par
Release-description mismatches | Bounded source discrepancies; paper impact unresolved | No claim that released code generated published outputs. No repair verified.\par
Existing steering issue \#44 | Known public concern; not reproduced | Original attribution preserved; target path and publication provenance unresolved.\par
Scientific decision / author adoption | Pending / not requested | No human-review declaration completed.\par
\section*{Supporting checks}
Original PDF, source-acquisition record, independent hash/title verification, extracted text, rendered source pages, pinned repository trees and per-file SHA-256 manifests are retained. Detailed inspection boundaries and line references are in claim\_audit.json and static\_code\_audit.json. All support files are under sources/0096 or verification/0096.\par
A guessed code path failed with ENOENT and was resolved by listing the actual repository tree.\par
No person is represented as having reviewed or approved the conclusions. Native compilation and current-PDF visual checks are recorded in this acceptance bundle.\par
\section*{Consulted references and evidence identifiers}
1. Adams E., Bai L., Lee M., Yu Y., AlQuraishi M. ICML 2025/PMLR 267:460-476. Exact publisher PDF identified by the source hash above. https://proceedings.mlr.press/v267/adams25a.html\par
2. InterProt source repository. Inspected commits fae2d3ae90220781289e27a26cb3b0826cf2b07e and 287c373817b06178603fc3a97bec95ebabb17ecc; acquired file lists and hashes retained. https://github.com/etowahadams/interprot\par
3. Public issue \#44, opened 4 December 2025. External allegation only; saved issue metadata preserves attribution. https://github.com/etowahadams/interprot/issues/44\par
4. Adams et al. bioRxiv 10.1101/2025.02.06.636901v2, 18 June 2025; metadata/version lookup only. https://doi.org/10.1101/2025.02.06.636901\par
5. Reviewer evidence: source\_identity\_verification.json, claim\_audit.json, static\_code\_audit.json, source\_lookup\_log.json and repository acquisition manifests, all within the 0096 source/verification directories.\par
\end{document}
